\documentclass[11pt,oneside]{report}
\usepackage[letterpaper,margin=1in,headheight=15pt]{geometry}
\usepackage{fontspec}
\setmainfont{texgyrepagella-regular.otf}[BoldFont=texgyrepagella-bold.otf,ItalicFont=texgyrepagella-italic.otf,BoldItalicFont=texgyrepagella-bolditalic.otf]
\setsansfont{texgyreheros-regular.otf}[BoldFont=texgyreheros-bold.otf,ItalicFont=texgyreheros-italic.otf]
\setmonofont{lmmono10-regular.otf}[Scale=0.85]
\usepackage{amsmath,amssymb,amsthm,booktabs,longtable,array,multirow,calc}
\usepackage{graphicx,xcolor,tikz,pgfplots}
\pgfplotsset{compat=1.18}
\usetikzlibrary{arrows.meta,positioning,calc,fit,shapes.geometric,patterns,decorations.pathreplacing}
% UIUC official brand colors, verified 2026-10-23.
% https://brand.illinois.edu/visual-identity/color
\definecolor{IlliniOrange}{HTML}{FF5F05}
\definecolor{IlliniBlue}{HTML}{13294B}
\definecolor{Storm}{HTML}{707372}
\definecolor{LightStorm}{HTML}{C6C7C6}
\definecolor{Industrial}{HTML}{1D58A7}
\definecolor{Patina}{HTML}{007E8E}
\tikzset{>=Stealth, every node/.style={font=\small,text=IlliniBlue}, box/.style={draw=IlliniBlue,rounded corners=2pt,align=center,inner sep=6pt,execute at begin node=\hyphenpenalty10000\relax}, orangebox/.style={box,draw=IlliniOrange,very thick}, edge/.style={->,thick,IlliniBlue}, accent/.style={->,very thick,IlliniOrange}}
\pgfplotsset{every axis/.append style={width=13.5cm,height=6.5cm,axis line style={IlliniBlue},tick label style={font=\small,text=IlliniBlue},label style={font=\small,text=IlliniBlue},legend style={font=\small,draw=none,fill=white},grid style={LightStorm},major grid style={LightStorm},axis lines=left,clip=false}, every axis plot/.append style={thick}}

\usepackage{titlesec,fancyhdr,enumitem,caption,float,placeins,fvextra}
\usepackage{microtype,xurl}
\usepackage{hyperref}
\hypersetup{colorlinks=true,linkcolor=IlliniBlue,urlcolor=Industrial,citecolor=IlliniBlue,pdftitle={BADM 534: Data Analytics for Marketing Applications},pdfauthor={Nathan Yang}}
\usepackage{bookmark}
\titleformat{\chapter}[display]{\normalfont\huge\bfseries\raggedright\hyphenpenalty=10000\color{IlliniBlue}}{\chaptertitlename\ \thechapter}{8pt}{\huge}
\titlespacing*{\chapter}{0pt}{-18pt}{20pt}
\titleformat{\section}{\Large\bfseries\raggedright\hyphenpenalty=10000\color{IlliniBlue}}{\thesection}{0.7em}{}
\titleformat{\subsection}{\large\bfseries\raggedright\hyphenpenalty=10000\color{IlliniBlue}}{\thesubsection}{0.7em}{}
\setlength{\parindent}{1.3em}
\setlength{\parskip}{3pt}
\linespread{1.06}
\setlist{itemsep=3pt,topsep=5pt}
\setcounter{tocdepth}{1}
\setcounter{secnumdepth}{2}
\pagestyle{fancy}
\fancyhf{}
\fancyhead[L]{\small\sffamily BADM 534}
\fancyhead[R]{\small\sffamily Lecture Notes}
\fancyfoot[C]{\thepage}
\renewcommand{\headrulewidth}{0.4pt}
\newcommand{\tightlist}{\setlength{\itemsep}{0pt}\setlength{\parskip}{0pt}}
\DefineVerbatimEnvironment{Highlighting}{Verbatim}{breaklines,fontsize=\small,commandchars=\\\{\}}
\newtheorem{assumption}{Assumption}[chapter]
\captionsetup{font=small,labelfont=bf,labelsep=period}
\setlength{\emergencystretch}{4em}
\providecommand{\pandocbounded}[1]{#1}

\clubpenalty=10000
\widowpenalty=10000

\begin{document}
\hypersetup{pdftitle={Marketing Dialectics}}
\setcounter{chapter}{8}
\hypertarget{M09-marketing-dialectics}{%
\chapter{Marketing Dialectics}\label{M09-marketing-dialectics}}
\noindent{\large By Nathan Yang}\par\medskip

\noindent \Needspace{8\baselineskip}

\hypertarget{M09-when-a-commercial-result-becomes-a-claim-about-customers}{%
\section{When a Commercial Result Becomes a Claim About
Customers}\label{M09-when-a-commercial-result-becomes-a-claim-about-customers}}

\noindent \noindent A subscription manager must decide whether to retain a
cancellation sequence that raises paid renewal. One explanation is that
the sequence helps customers discover a service feature they value.
Another is that it makes leaving sufficiently difficult that some
customers continue paying against their wishes. Both mechanisms can
produce the same billing record. The decision therefore depends on the
reasons for renewal, the customer's understanding of the offer, and the
practical ability to refuse it.

The claim that higher renewal demonstrates customer benefit contains a
premise: continued payment expresses an informed wish to continue at the
stated terms. That premise requires evidence. A manager can calculate
renewal accurately and still draw an unjustified conclusion about the
customer. Equally, the existence of a retention offer does not establish
manipulation. A relevant discount, clearly explained and freely refused,
can improve the customer's available options.

Kant's critique of dialectical reasoning asks how reason comes to treat
its own demands for explanation as knowledge about objects. In a
marketing application, the corresponding discipline is to identify what
connects an observed result to a broader claim. This is an adaptation of
Kant's critical approach to an empirical decision. A disputed
interpretation of a customer metric is not, by itself, a transcendental
antinomy. The philosophical connection concerns the examination of
premises and the limits of a conclusion.

Three questions organize the decision. What does the evidence establish
about the practice? Which assumptions connect that finding to customer
interests? Which ethical principles justify the action? A favorable
answer to one question leaves the others substantive. An experiment may
identify an increase in purchases while leaving the quality of consent
unresolved. A transparent offer may respect a customer's choice while
producing too little revenue to remain commercially viable.

\begin{figure}[htbp]
\centering
\begin{tikzpicture}[node distance=1.05cm and 1.05cm]
\node[box,text width=3cm] (practice) at (0,0) {Marketing practice};
\node[box,text width=3.5cm] (sales) at (5.2,1.5) {Commercial effect\\Sales, margin, renewal};
\node[orangebox,text width=3.5cm] (customer) at (5.2,-1.5) {Customer outcome\\Understanding, value, exit};
\draw[edge] (practice.east) -- (sales.west);
\draw[accent] (practice.east) -- (customer.west);
\node[align=center,font=\footnotesize] at (1.6,1.5) {Causal\\design};
\node[align=center,font=\footnotesize] at (1.6,-1.7) {Measure and\\comparison};
\node[box,text width=3.5cm] (judgment) at (10.5,0) {Ethical judgment\\Stated principle of action};
\draw[edge] (sales.east) -- (judgment.north west);
\draw[accent] (customer.east) -- (judgment.south west);
\node[align=center,text width=11cm] at (5.2,-3.1) {A measured sales effect leaves the customer measure and ethical premise substantive.};
\end{tikzpicture}

\caption{A commercial outcome, a customer outcome, and an ethical judgment require different supporting premises.}\label{fig:M09_claim_scope}
\end{figure}

The upper path connects an intervention to a measured commercial outcome
through a causal design. A second path requires a valid customer outcome
measure and a defined comparison. The ethical judgment additionally
depends on a stated principle concerning treatment of persons. No arrow
permits sales to stand in for all three conclusions.

\Needspace{8\baselineskip}

\hypertarget{M09-kants-dialectic-and-the-limits-of-reason}{%
\section{Kant's Dialectic and the Limits of
Reason}\label{M09-kants-dialectic-and-the-limits-of-reason}}

\noindent \hypertarget{M09-the-search-for-an-unconditioned-explanation}{%
\subsection{The Search for an Unconditioned
Explanation}\label{M09-the-search-for-an-unconditioned-explanation}}

\noindent \noindent In the \emph{Critique of Pure Reason}, the understanding
organizes experience through concepts and judgments. Reason seeks
relations among those judgments and asks for more complete explanations.
Given something conditioned, it looks for the condition on which it
depends. Repeating that question creates a demand for an unconditioned
ground that would complete the explanation. The difficulty arises when
the demand for completeness is treated as establishing an object that
theoretical reason can know independently of possible experience. Kant
develops this account in the introduction to the Transcendental
Dialectic, especially A293/B349 through A308/B365. The A and B
references identify the original first- and second-edition pagination.
\href{https://www.gutenberg.org/files/4280/4280-h/4280-h.htm}{Kant,
\emph{Critique of Pure Reason}}

Kant distinguishes transcendental illusion from a mistake that
disappears once a faulty inference is corrected. The attraction of an
unconditional explanation persists after criticism has exposed its
source. The task of critique is to prevent that attraction from becoming
an unjustified judgment. This distinction gives an important reason to
revisit assumptions. Recognizing a tempting inference once does not
guarantee that an attractive new result will cease to invite it.

The three major divisions of the Dialectic concern the soul, the world
as a whole, and God. The paralogisms examine inferences from the unity
of thinking to claims about the soul as an object. The antinomies
concern opposed arguments about the world considered as a completed
totality. The ideal of pure reason concerns the supposed highest being
and attempts to establish its existence. These are specific arguments
within Kant's critical philosophy. Their marketing relevance comes from
the practice of examining the transition from a valid use of reason to
an unsupported claim, rather than from equating a customer segment,
dataset, or company with one of these metaphysical objects.

\hypertarget{M09-antinomies-and-the-shared-premise}{%
\subsection{Antinomies and the Shared
Premise}\label{M09-antinomies-and-the-shared-premise}}

\noindent \noindent Kant presents four cosmological conflicts: whether the world
has limits in time and space; whether composite substances ultimately
consist of simple parts; whether causality through freedom is necessary
alongside natural causality; and whether an absolutely necessary being
exists. The apparent strength of opposing arguments calls attention to
their common framework. In Kant's diagnosis, treating the world of
appearances as a completed whole of things in themselves generates the
conflict. His resolutions depend on the particular antinomy. The
mathematical and dynamical antinomies do not all receive an identical
verdict. The \emph{Prolegomena}, §§50--54, supplies a shorter treatment.
\href{https://www.gutenberg.org/files/52821/52821-h/52821-h.htm}{Kant,
\emph{Prolegomena}}

The third antinomy is especially relevant to discussions of prediction
and agency. Its thesis argues for a causality through freedom; its
antithesis argues that events occur through natural laws. Kant's
critical solution distinguishes the explanation of events as appearances
from the intelligible standpoint under which freedom can be thought. The
first Critique's resolution does not supply an empirical demonstration
that a particular choice was free. Interpretations of transcendental
idealism and the relationship between these standpoints remain disputed.
\href{https://plato.stanford.edu/entries/kant-transcendental-idealism/}{Stang,
\emph{Kant's Transcendental Idealism}}

A marketing model that predicts a purchase therefore does not settle a
metaphysical question about freedom. Prediction can reflect stable
preferences, recurring needs, constraints, exposure, or a combination of
these. Conversely, an unpredictable choice can arise from poor
measurement or random error. Neither predictive success nor prediction
failure establishes whether a firm has respected a person's capacity to
judge and choose.

\begin{figure}[htbp]
\centering
\begin{tikzpicture}
\node[box,text width=4.6cm] (a) at (0,0) {Higher clicks prove a better fit\\with customers' goals.};
\node[box,text width=4.6cm] (b) at (7,0) {Changed clicks prove that\\customers' goals were overridden.};
\node[orangebox,text width=8cm] (p) at (3.5,-2.1) {Shared premise: recorded clicks fully express the person's considered ends.};
\draw[edge] (a.south) -- (p.north west);
\draw[edge] (b.south) -- (p.north east);
\node[box,text width=9cm] (q) at (3.5,-4.3) {Examine information, the customer's purpose, and the ability to refuse or revise the recommendation.};
\draw[accent] (p) -- (q);
\end{tikzpicture}

\caption{Opposing claims can inherit the same unsupported premise; examining that premise can change the question.}\label{fig:M09_shared_premise}
\end{figure}

The marketing example places two claims above a shared assumption:
observed clicks fully express customers' considered ends. One claim
treats higher clicks as proof of a better fit with customers' goals; the
other treats changed clicks as proof that those goals were overridden.
Examining the shared assumption allows a more specific evaluation of
information, goals, and control. This argument map is a contemporary
analogy. It is not one of Kant's four cosmological antinomies.

\hypertarget{M09-regulative-ideas-and-productive-inquiry}{%
\subsection{Regulative Ideas and Productive
Inquiry}\label{M09-regulative-ideas-and-productive-inquiry}}

\noindent \noindent The appendix to the Dialectic gives reason a constructive
role. Ideas can guide the search for systematic unity without providing
constitutive knowledge of the objects they represent. A research program
can seek more connected explanations while preserving the conditions
under which each claim is warranted. Kant distinguishes this regulative
use from treating an idea as if it supplied a determinate object of
knowledge. The interpretation of the status of these ideas is itself
debated; Gava distinguishes different kinds of regulative assumptions
rather than treating every idea as a known fiction.
\href{https://doi.org/10.1515/opphil-2022-0202}{Gava (2022)}

In a marketing adaptation, a commitment to understanding customer
benefit can guide the collection of information about learning, regret,
comprehension, and the ability to leave. It does not justify announcing
that benefit has already been measured by a single score. A useful ideal
generates better questions. It becomes misleading when its desirable
name is attached to a convenient outcome and then treated as a finding.

Two further uses of dialectic broaden the reading. In the \emph{Critique
of Practical Reason}, Book II, Kant examines the highest good and the
relation between virtue and happiness. His argument supplies no
empirical guarantee that ethical firms earn the highest profits. In the
\emph{Critique of Judgement}, §§55--57, the antinomy of taste concerns
the basis of judgments that claim assent without being demonstrable
through determinate concepts. That problem differs from estimating which
advertisement a customer prefers. These distinctions preserve the
specific object of each argument.
\href{https://www.gutenberg.org/files/5683/5683-h/5683-h.htm}{Kant,
\emph{Practical Reason}};
\href{https://www.gutenberg.org/files/48433/48433-h/48433-h.htm}{Kant,
\emph{Judgement}}

\Needspace{8\baselineskip}

\hypertarget{M09-a-critical-examination-of-marketing-claims}{%
\section{A Critical Examination of Marketing
Claims}\label{M09-a-critical-examination-of-marketing-claims}}

\noindent \noindent Consider the claim that personalized offers improve customer
welfare. Begin by stating the practice: which information enters the
rule, which offers can be shown, whether prices change, and what
alternatives remain available. Next identify the comparison. The claim
could concern personalization relative to a common offer, a
customer-selected rule, or a previous algorithm. These comparisons can
have different effects. A broad phrase such as personalization is
beneficial leaves both treatment and benchmark unspecified.

A descriptive claim records an observed pattern. A predictive claim
concerns performance on a specified population and information set. A
causal claim concerns what changes under an intervention relative to an
alternative. A normative claim concerns what should be permitted or
chosen. The evidence must fit the claim. An accurate description of
sales among exposed customers does not by itself identify incremental
sales; an incremental-sales estimate does not by itself determine
whether the intervention is ethically acceptable.

The most useful objection identifies a premise on which the
recommendation depends. For an advertising claim, that premise might be
that exposed and unexposed customers would have bought at the same rate
without exposure. For a cancellation claim, it might be that continued
billing expresses an informed wish to renew. For a recommendation claim,
it might be that the platform's engagement objective adequately
represents the customer's purpose. Once the premise is stated, the
analyst can ask which observation, design, or argument would support it
and which would undermine it.

Disagreement then has several possible resolutions. A claim can be
accepted within a specified population and horizon. Its scope can be
narrowed, its mechanism revised, or its recommendation rejected.
Judgment can also be suspended when a consequential premise remains
unresolved. There is no requirement to compromise between a supported
claim and an unsupported objection. Kant's dialectic should not be
presented as a mechanical sequence in which every thesis requires an
antithesis and an inevitable synthesis. Later dialectical traditions
require their own interpretation.
\href{https://plato.stanford.edu/entries/hegel-dialectics/}{Maybee,
\emph{Hegel's Dialectics}}

Critical examination also applies to the criticism. A claim that all
advertising overrides agency must account for advertising that
communicates useful price, availability, or quality information. A claim
that personalization always benefits customers must account for
inaccurate inferences, concealed incentives, and customers whose goals
differ from the model's objective. The burden is to explain the scope of
each argument and the evidence that bears on it.

\Needspace{8\baselineskip}

\hypertarget{M09-respect-for-persons-and-commercial-exchange}{%
\section{Respect for Persons and Commercial
Exchange}\label{M09-respect-for-persons-and-commercial-exchange}}

\noindent \hypertarget{M09-means-ends-and-the-customers-judgment}{%
\subsection{Means, Ends, and the Customer's
Judgment}\label{M09-means-ends-and-the-customers-judgment}}

\noindent \noindent Kant's ethical analysis adds a distinct question to the
critique of knowledge: how should one act toward rational persons? The
humanity formulation of the categorical imperative requires treating
humanity as an end and never merely as a means. A normal exchange uses
the services of another person in pursuing an end. Its moral assessment
depends on whether that person is also respected as capable of forming
and pursuing ends. The qualifier merely therefore matters. Commercial
activity is not condemned simply because a seller seeks revenue.
\href{https://www.gutenberg.org/files/5682/5682-h/5682-h.htm}{Kant,
\emph{Groundwork}, second section, Akademie 4:429}

A customer who knowingly chooses a useful service and can revise that
choice may benefit alongside the seller. A process that relies on
concealed terms or prevents an expressed refusal presents a different
relation between the parties. The concern is that the customer's
reasoning and choice become obstacles to be circumvented. This
evaluation requires the details of the practice: what the firm
communicates, what it conceals, what the customer understands, and what
happens when the customer says no.

Santori's interpretation of Kant develops the possibility of market
exchange organized around mutual assistance and the spirit of mutually
beneficial contracts. It provides a useful counterweight to the
assumption that a profit motive makes exchange intrinsically
objectionable. The argument is philosophical; it does not establish that
observed contracts are informed, fair, or beneficial in every setting.
\href{https://link.springer.com/article/10.1007/s40926-024-00252-z}{Santori
(2024)}

\hypertarget{M09-universalization-and-the-maxim-of-a-practice}{%
\subsection{Universalization and the Maxim of a
Practice}\label{M09-universalization-and-the-maxim-of-a-practice}}

\noindent \noindent A maxim specifies a circumstance, an action, and a purpose.
For example: when a customer seeks to leave, conceal the effective
cancellation route so that the customer continues paying. Stating the
maxim exposes the role of concealment in the business objective. Kant's
universal-law formulation asks whether a maxim can be willed as a
universal law, with distinctions between contradictions in conception
and in willing. This is a test of a principle of action, rather than a
vote on whether the practice is popular. {[}Kant, \emph{Groundwork},
second section, Akademie 4:421{]}

The moral difficulty becomes especially clear when a promise of easy
cancellation is used to obtain enrollment while the firm intends to
frustrate cancellation later. The strategy depends on the customer
relying on an assurance the firm does not intend to honor. A truthful
and accessible renewal reminder has a different maxim. Its function can
be to help customers act on their plans while making an optional offer.
The formulation of the maxim must reflect the actual practice and
purpose; redescribing concealment as customer service does not answer
the objection.

Kantian autonomy concerns rational self-legislation. A recorded click is
a behavioral event, and a preference survey is a measurement instrument.
Neither is a complete operational definition of autonomy. Comprehension,
refusal, revision, and effective exit are useful empirical indicators
for evaluating a marketing practice. They leave the philosophical
meaning of autonomy broader than any single observed response.
\href{https://plato.stanford.edu/entries/kant-moral/}{Johnson and
Cureton, \emph{Kant's Moral Philosophy}}

\hypertarget{M09-ethical-reasons-professional-standards-and-law}{%
\subsection{Ethical Reasons, Professional Standards, and
Law}\label{M09-ethical-reasons-professional-standards-and-law}}

\noindent \noindent Laczniak and Murphy connect marketing ethics to duties toward
customers, other stakeholders, and the broader marketing system. Their
account encourages attention to effects that a narrow transaction
measure omits. The American Marketing Association's current ethics
statement likewise articulates professional values, including honesty,
responsibility, equity, transparency, and citizenship. It is an
aspirational professional statement. A professional code, a
philosophical argument, and an enforceable legal obligation have
different sources of authority.
\href{https://epublications.marquette.edu/market_fac/45/}{Laczniak and
Murphy (2006)}; \href{https://www.ama.org/ama-statement-of-ethics/}{AMA,
\emph{Statement of Ethics}}

Customer benefit, respectful treatment, and legal compliance should each
receive a reasoned evaluation. A lawful practice can still warrant
ethical criticism. A beneficial effect does not erase deception. An
allegation of deception requires evidence about the communication and
its likely meaning. In a cancellation case, the applicable jurisdiction,
transaction, and current rule must be identified before making a legal
claim. The
\href{../01_Resources/13_Consumer_Agency_and_Recent_Research/Policy/README.md}{course
policy source record} separates the status of federal rulemaking from
other applicable cancellation requirements.

\Needspace{8\baselineskip}

\hypertarget{M09-worked-case-identical-renewal-different-customer-intent}{%
\section{Worked Case: Identical Renewal, Different Customer
Intent}\label{M09-worked-case-identical-renewal-different-customer-intent}}

\noindent \hypertarget{M09-the-population-and-the-missing-distinction}{%
\subsection{The Population and the Missing
Distinction}\label{M09-the-population-and-the-missing-distinction}}

\noindent \noindent Consider two entirely constructed populations, A and B, each
containing 1,000 customers eligible for one monthly renewal. Each
customer appears once. Billing is observed during the same renewal
month. Whether the customer wanted to renew at the stated terms is
stipulated before billing and independently of whether a charge
occurred. No actual platform records, randomized treatment, or estimated
preferences enter this example.

In A, 880 customers wanted renewal and were billed; 20 did not want
renewal and were billed. In B, these counts are 600 and 300. Both
populations also contain 20 customers who wanted renewal but were not
billed and 80 who neither wanted renewal nor were billed. Each
population therefore has 900 billed customers and a billed-renewal rate
of 90\%. Unwanted billing represents 2\% of all customers in A and 30\%
in B.

\begin{table}[H]
\centering\small
\begin{tabular}{@{}lrr@{}}
\toprule
Customer State & Population A & Population B \\
\midrule

Wanted renewal and billed & 880 & 600 \\
Did not want renewal and billed & 20 & 300 \\
Wanted renewal and not billed & 20 & 20 \\
Did not want renewal and not billed & 80 & 80 \\
All eligible customers & 1,000 & 1,000 \\
\bottomrule
\end{tabular}
\end{table}

To distinguish the billing outcome from its composition, let \(B_i\)
equal one when customer \(i\) is billed and let \(W_i\) equal one when
that customer wanted renewal before billing. For each population of size
\(N\), define

\[
R=\frac{1}{N}\sum_{i=1}^{N}B_i,
\qquad
U=\frac{1}{N}\sum_{i=1}^{N}B_i(1-W_i).
\]

Both quantities are customer shares over the renewal month. \(R\) is
billed renewal; \(U\) is unwanted billing among all eligible customers.
Holding \(R\) fixed leaves \(U\) undetermined because the billing total
does not reveal the intent of those billed. Among billed customers, the
unwanted shares instead use 900 as the denominator, giving approximately
2.22\% in A and 33.33\% in B. A conditional share and a population share
answer different questions.

\begin{figure}[htbp]
\centering
\begin{tikzpicture}[x=1cm,y=1cm]
\node[anchor=east] at (-.3,0.4) {A};
\node[anchor=east] at (-.3,-1.3) {B};
\fill[IlliniBlue] (0,0) rectangle (8.8,.8);
\fill[IlliniOrange] (8.8,0) rectangle (9,.8);
\fill[LightStorm] (9,0) rectangle (10,.8);
\fill[IlliniBlue] (0,-1.7) rectangle (6,-.9);
\fill[IlliniOrange] (6,-1.7) rectangle (9,-.9);
\fill[LightStorm] (9,-1.7) rectangle (10,-.9);
\node[text=white] at (4.4,.4) {880 wanted and billed};
\node[text=white] at (3,-1.3) {600 wanted and billed};
\node[font=\footnotesize] at (7.5,-1.3) {300 unwanted};
\node[anchor=south,align=center,font=\footnotesize] (label) at (8.0,1.4) {20 unwanted};
\draw[accent] (label.south) -- (8.9,.82);
\node[anchor=west,align=left,font=\footnotesize] at (10.2,.4) {100 not\\billed};
\node[anchor=west,align=left,font=\footnotesize] at (10.2,-1.3) {100 not\\billed};
\draw[IlliniBlue] (0,-2.25) -- (10,-2.25);
\foreach \x/\t in {0/0,5/50,10/100} {
 \draw[IlliniBlue] (\x,-2.25) -- (\x,-2.4);
 \node[anchor=north] at (\x,-2.4) {\t\%};
}
\node[align=center] at (5,-3.4) {Share of all 1,000 customers in each synthetic population\\Billed renewal: 90\% in both; unwanted billing: 2\% versus 30\%};
\end{tikzpicture}

\caption{Two synthetic populations have 90\% billed renewal; unwanted billing occupies 2\% of all customers in A and 30\% in B.}\label{fig:M09_retention_intent}
\end{figure}

The bars divide each full population into wanted and billed, unwanted
and billed, and not billed. Their common billing total conceals a
different composition. These are observationally equivalent populations
with respect to the billing-rate statistic. The construction supplies no
estimate of the effect of replacing one cancellation interface with
another. It also does not assign a monetary welfare loss to unwanted
billing.

\hypertarget{M09-reproducible-calculation-and-missing-intent}{%
\subsection{Reproducible Calculation and Missing
Intent}\label{M09-reproducible-calculation-and-missing-intent}}

\noindent \noindent The program below constructs every customer record from the
eight specified cells, exports the synthetic data, and calculates the
relevant counts and shares. It retains non-billed customers in the
population denominator. The summary function permits unknown intent and
reports bounds: the lower bound counts observed unwanted billing, while
the upper bound also counts every billed customer whose intent is
unknown. Correctly measured known labels are assumed for those bounds.
Missing intent supplies no evidence of satisfaction.

\FloatBarrier
\Needspace{8\baselineskip}
Python source:
\href{../03_Data_and_Code/Marketing_Dialectics/evaluate_retention.py}{evaluate\_retention.py}

\begin{Verbatim}[breaklines,breakanywhere,fontsize=\small]
"""Compare two constructed customer populations with identical billed renewal.

Python 3.10+; standard library only. No network calls or real customer records.
Run: python3 03_Data_and_Code/Marketing_Dialectics/evaluate_retention.py
"""
import csv
import json
from pathlib import Path

ROOT = Path(__file__).resolve().parent


def construct_customers():
    # Each tuple is (billed this month, wanted renewal before billing, count).
    # Counts are imposed teaching inputs, not estimated responses to a policy.
    cells = {
        "A": [(True, True, 880), (True, False, 20),
              (False, True, 20), (False, False, 80)],
        "B": [(True, True, 600), (True, False, 300),
              (False, True, 20), (False, False, 80)],
    }
    rows = []
    for population, groups in cells.items():
        number = 0
        for billed, wanted, count in groups:
            for _ in range(count):
                number += 1
                rows.append({"population": population,
                             "customer_id": f"{population}{number:04d}",
                             "billed": billed, "wanted_renewal": wanted})
    return rows


def summarize(rows):
    if not rows:
        raise ValueError("The population must contain customers.")
    keys = [(r["population"], r["customer_id"]) for r in rows]
    if len(set(keys)) != len(keys):
        raise ValueError("Duplicate customer within population.")
    if len({r["population"] for r in rows}) != 1:
        raise ValueError("Summarize one population at a time.")
    for r in rows:
        if type(r["billed"]) is not bool:
            raise ValueError("Billing must be observed as True or False.")
        if r["wanted_renewal"] is not None and type(r["wanted_renewal"]) is not bool:
            raise ValueError("Intent must be True, False, or None.")
    n = len(rows)
    billed = sum(r["billed"] for r in rows)
    wanted_billed = sum(r["billed"] and r["wanted_renewal"] is True for r in rows)
    unwanted_billed = sum(r["billed"] and r["wanted_renewal"] is False for r in rows)
    unknown_billed = sum(r["billed"] and r["wanted_renewal"] is None for r in rows)
    missed_wanted = sum(not r["billed"] and r["wanted_renewal"] is True for r in rows)
    assert wanted_billed + unwanted_billed + unknown_billed == billed
    return {
        "customers": n, "billed": billed, "billed_rate": billed / n,
        "wanted_billed": wanted_billed, "unwanted_billed": unwanted_billed,
        "unknown_intent_billed": unknown_billed,
        "wanted_but_not_billed": missed_wanted,
        "unwanted_billed_lower_share_all": unwanted_billed / n,
        "unwanted_billed_upper_share_all": (unwanted_billed + unknown_billed) / n,
        "unwanted_share_among_billed_lower": unwanted_billed / billed if billed else None,
    }


def main():
    rows = construct_customers()
    with (ROOT / "synthetic_customer_months.csv").open("w", newline="") as f:
        writer = csv.DictWriter(f, fieldnames=list(rows[0]))
        writer.writeheader()
        writer.writerows(rows)
    results = {p: summarize([r for r in rows if r["population"] == p])
               for p in sorted({r["population"] for r in rows})}
    output = {
        "data_type": "Entirely synthetic; two constructed populations",
        "unit": "One customer over one renewal month",
        "intent_timing": "Before billing; stipulated independently of billing",
        "comparison": "Descriptive identification example, not a policy effect",
        "results": results,
    }
    (ROOT / "results.json").write_text(json.dumps(output, indent=2) + "\n")
    print(json.dumps(results, indent=2))


if __name__ == "__main__":
    main()
\end{Verbatim}

The output reports 900 billed customers in each population, with 20 and
300 unwanted charges respectively. Both contain 20 customers who wanted
renewal but were not billed. The five independent tests in the example
folder also check a four-customer missing-intent case, a population with
no billing, duplicate records, and string values that could be mistaken
for booleans.

\hypertarget{M09-from-the-constructed-case-to-a-decision}{%
\subsection{From the Constructed Case to a
Decision}\label{M09-from-the-constructed-case-to-a-decision}}

\noindent \noindent In a real service, intent would need a defined measurement
procedure. Asking after a charge can capture regret, newly learned
information, or a response to the charge itself. Asking before renewal
can affect the subsequent choice. Nonrespondents can differ from
respondents. The timing and wording of the question therefore belong in
the interpretation, along with response rates and the treatment of
unknown answers.

A useful design compares a clearly disclosed, direct renewal or
cancellation path with a proposed alternative and measures more than the
payment outcome. Relevant outcomes include understood terms, successful
completion of the customer's requested action, unwanted subsequent
charges, and later evaluation of the service. Commercial viability
remains part of the decision. These measurements make it possible to
evaluate whether a useful offer improves an informed choice or whether
extra revenue depends on frustrating that choice.

Luguri and Strahilevitz experimentally varied interfaces for declining
an identity-theft protection plan and found that the interface affected
continued enrollment. Their study motivates attention to the path
through which a recorded decision is obtained. Mathur and coauthors
instead document interface practices across shopping websites; their
crawl supplies descriptive evidence about prevalence in that sample. The
FTC's 2022 staff report provides additional examples and a regulatory
taxonomy. These sources answer different questions about effects,
prevalence, and regulatory concern.
\href{https://oag.ca.gov/sites/all/files/agweb/pdfs/privacy/shining-a-light-on-dark-patterns.pdf}{Luguri
and Strahilevitz, 2019 working paper};
\href{https://arxiv.org/abs/1907.07032}{Mathur et al.~(2019)};
\href{https://www.ftc.gov/reports/bringing-dark-patterns-light}{FTC
(2022)}

\Needspace{8\baselineskip}

\hypertarget{M09-advertising-effectiveness-and-the-meaning-of-a-successful-test}{%
\section{Advertising Effectiveness and the Meaning of a Successful
Test}\label{M09-advertising-effectiveness-and-the-meaning-of-a-successful-test}}

\noindent \noindent A retailer must decide whether an advertising campaign
deserves additional spending. Customers exposed to the ads purchase more
frequently than customers who were not exposed. The analyst proposes
that this difference measures the campaign's effect. A competing
explanation is that the targeting system selected customers already more
likely to purchase. The dispute concerns a counterfactual premise: what
the exposed customers would have done without the campaign.

Gordon and coauthors compare observational estimates with randomized
advertising experiments conducted at Facebook. Their results show that
extensive observed information does not reliably remove the discrepancy
in those campaigns. The study supports careful identification in
advertising measurement; it does not establish that every observational
design is invalid. Google Ads' Conversion Lift documentation describes
an industry approach to comparing groups with and without advertising
exposure. Provider documentation explains an implemented service, while
a particular evaluation still requires attention to its assignment,
eligibility, measurement, and study type.
\href{https://www.kellogg.northwestern.edu/faculty/gordon_b/files/fb_comparison.pdf}{Gordon
et al., 2018 working paper, later published 2019};
\href{https://support.google.com/google-ads/answer/12003020?hl=en}{Google
Ads, \emph{About Conversion Lift}}

Suppose a well-executed randomized study establishes positive
incremental sales for the eligible customers and observation window. The
appropriate next question depends on the managerial claim. A budget
decision requires contribution after relevant costs and consideration of
how effects change at greater spending. A customer-benefit claim
requires evidence about the product's consequences and the way the offer
was communicated. A successful sales experiment does not answer these
additional questions simply by having a credible assignment mechanism.

A Kantian evaluation therefore preserves the experimental finding while
examining the principle of the practice. A truthful advertisement that
helps customers discover a useful option can serve both parties' ends.
An advertisement that secures the same incremental sales through
concealed material terms presents a different ethical case. The
objection concerns the means of obtaining the result. It can remain
decisive even when the sales effect is precisely measured.

\Needspace{8\baselineskip}

\hypertarget{M09-personalization-learning-and-the-ability-to-revise-a-goal}{%
\section{Personalization, Learning, and the Ability to Revise a
Goal}\label{M09-personalization-learning-and-the-ability-to-revise-a-goal}}

\noindent \noindent A recommendation system can reduce search costs by selecting
items that fit a customer's interests. It can also shape the information
from which later interests are inferred. Netflix's account of artwork
personalization describes selecting different images to help members
discover titles. The description supplies a concrete industry mechanism:
presentation changes which features of an item attract attention. It
does not establish that every resulting viewing decision serves the
member's considered purpose.
\href{https://netflixtechblog.com/artwork-personalization-c589f074ad76}{Chandrashekar,
Amat, Basilico, and Jebara (2017)}

Two claims deserve scrutiny. One says that selecting an item reveals
benefit because the customer clicked. The other says that an algorithm's
influence on attention necessarily removes agency. Both can ignore the
process by which the person understands and evaluates the
recommendation. Useful information can change a preference through
learning. A misleading representation can induce a choice the customer
would reject after learning the relevant facts. The substantive
difference concerns the information and the person's capacity to
reflect, rather than whether behavior changed at all.

\begin{figure}[htbp]
\centering
\begin{tikzpicture}
\node[box,text width=2.8cm] (rank) at (0,0) {Ranking rule};
\node[box,text width=2.8cm] (exposure) at (4.7,0) {Exposure};
\node[box,text width=2.8cm] (action) at (9.4,0) {Recorded behavior};
\draw[edge] (rank) -- (exposure);
\draw[edge] (exposure) -- (action);
\draw[edge] (action.south) -- ++(0,-1.5) -| node[pos=.25,below]{Training feedback} (rank.south);
\node[orangebox,text width=6.2cm] (goals) at (4.7,2.4) {Customer purpose, informed refusal, and goal revision};
\draw[accent] (goals.west) -| (rank.north);
\draw[accent] (goals.east) -| (action.north);
\node[align=center,text width=11.3cm] at (4.7,-3.2) {Behavior is recorded under the current exposure rule.\\Customer purposes require additional information.};
\end{tikzpicture}

\caption{A recommendation affects exposure and recorded behavior; customer goals and the ability to revise the rule supply information outside that feedback loop.}\label{fig:M09_goal_revision}
\end{figure}

The closed loop runs from the ranking rule through exposure and behavior
back to the training data. An additional path lets the customer state or
revise a goal and refuse a recommendation. That path does not prove
autonomy by itself. It identifies decisions and observations that an
engagement-only objective omits.

The recurring mHealth case makes the distinction concrete. A
recommendation based on past food logging might help someone organize
meals, maintain variety, or pursue a stated nutritional objective.
Frequent logging can also reflect the design of reminders and the burden
of recording. The logs record entries, rather than all consumption or
independently verified health outcomes. A recommendation should
therefore specify the customer's goal, the source and validity of
nutritional information, and the ability to change or decline the goal.
Health-related claims require appropriate domain evidence; engagement
cannot supply it.

An intervention can preserve customer control by making the
recommendation's purpose intelligible, avoiding materially misleading
descriptions, and allowing users to choose another objective or leave
the recommendation process. These are concrete design proposals to
evaluate. Their effects on understanding, use, and customer outcomes
must still be measured. A philosophical commitment gives direction to
that evaluation without predetermining its empirical result.

\Needspace{8\baselineskip}

\hypertarget{M09-environmental-claims-context-and-the-scope-of-an-inference}{%
\section{Environmental Claims, Context, and the Scope of an
Inference}\label{M09-environmental-claims-context-and-the-scope-of-an-inference}}

\noindent \noindent An energy company can truthfully describe a lower-carbon
product while creating a misleading impression about the balance of its
overall activities. The relevant object changes from a particular
offering to the company as a whole. Customers may reasonably interpret
the combination of imagery, brand references, and qualifying information
more broadly than an isolated sentence. Assessing the claim therefore
requires the communication in context.

In its June 7, 2023 Shell ruling, the UK Advertising Standards Authority
upheld the complaint concerning the overall environmental impression of
several ads. It found that information about the balance of the
company's activities was material to the impression conveyed. A separate
issue concerning the renewable-electricity claim was not upheld. The
decision consequently supports a specific conclusion about the evaluated
communications and omissions. It is not a finding that every statement
about the company's lower-carbon products was false.
\href{https://www.asa.org.uk/rulings/shell-uk-ltd-g22-1170842-shell-uk-ltd.html}{ASA,
Shell UK Ltd, June 7, 2023}

The ASA reached a different outcome on April 9, 2025 for another Shell
advertisement. That ad included higher- and lower-carbon activities and
qualifying information about the investment mix. The ASA found that the
overall environmental impression was not misleading and did not uphold
the complaint. Its assessment applied the code wording relevant to the
earlier complaint, as the ruling explains. Comparing the cases requires
attention to the actual advertisements and applicable assessment, rather
than treating the company name as a verdict.
\href{https://www.asa.org.uk/rulings/shell-uk-ltd-g24-1248246-shell-uk-ltd.html}{ASA,
Shell UK Ltd, April 9, 2025}

\begin{figure}[htbp]
\centering
\begin{tikzpicture}
\node[box,text width=3cm] (product) at (0,0) {Product claim\\A particular offering};
\node[orangebox,text width=3cm] (portfolio) at (5,0) {Portfolio claim\\The balance of activities};
\node[box,text width=3cm] (future) at (10,0) {Transition claim\\A future trajectory};
\draw[edge,dashed] (product) -- node[above,align=center,font=\footnotesize]{Requires\\context} (portfolio);
\draw[edge,dashed] (portfolio) -- node[above,align=center,font=\footnotesize]{Requires\\evidence} (future);
\node[align=center,text width=3cm] at (0,-2) {Product definition\\Relevant comparison};
\node[align=center,text width=3cm] at (5,-2) {Share and\\denominator\\Material omissions};
\node[align=center,text width=3cm] at (10,-2) {Time horizon\\Implemented commitments};
\node[align=center,text width=12cm] at (5,-3.5) {A supported narrow statement can invite an unsupported broader inference.};
\end{tikzpicture}

\caption{A claim about a product does not by itself establish a claim about a portfolio or a future transition.}\label{fig:M09_environmental_scope}
\end{figure}

The three boxes distinguish the product, the business portfolio, and the
future plan. A broader conclusion requires information about the
relevant share, time horizon, and basis of comparison. The decision is
to communicate the supportable claim with material context, or to change
the claim when the broader impression exceeds its support.

The regulatory findings and the philosophical evaluation remain
distinct. The ASA rulings apply UK advertising standards to particular
communications. A Kantian objection to deception concerns respect for
the customer's judgment. The same advertisement can be examined through
both lenses, but one source of authority should not be silently
substituted for the other. A comprehensive ethical judgment about a
company's environmental conduct would require additional evidence and
additional premises.

\Needspace{8\baselineskip}

\hypertarget{M09-ai-assisted-criticism-and-independent-judgment}{%
\section{AI-Assisted Criticism and Independent
Judgment}\label{M09-ai-assisted-criticism-and-independent-judgment}}

\noindent \noindent An LLM can help identify missing premises and plausible
alternative explanations. It can also produce an articulate defense of a
conclusion that the user has already selected. Fluency therefore
increases the importance of checking whether an objection is relevant,
whether its source exists, and whether the source supports the
proposition attributed to it. Agreement between two generated responses
provides little assurance when both repeat the same premise or
inaccurate citation.

A useful prompt makes the object of criticism explicit:

\begin{quote}
Evaluate the claim that this renewal intervention benefits customers.
State the population, intervention, comparison, outcome, and horizon.
Separate descriptive, causal, and ethical premises. Identify the
strongest plausible alternative explanation and the observation that
would distinguish it. Evaluate the treatment of customer understanding,
refusal, and revision. Cite accessible primary sources for factual
claims and identify unverified statements. Recommend acceptance,
restriction, redesign, rejection, or further evidence, with reasons.
\end{quote}

The analyst should first check the reconstructed argument against the
actual proposal. A system can criticize a practice the firm never
proposed or overlook a restriction that changes the conclusion. Next,
verify sources and calculations independently. Finally, state a decision
in the analyst's own terms and identify the evidence that would change
it. This process can support critical thought while leaving
responsibility for the judgment with the person making the
recommendation.

Kant's essay on enlightenment connects maturity with the use of one's
own understanding and the freedom to examine claims publicly. It offers
a useful reason to resist treating a generated answer as an
unquestionable authority. Independent judgment can include careful
reliance on expertise; it requires understanding why the expertise is
relevant and how its conclusion is supported. The essay's distinction
between public and private uses of reason belongs to its
eighteenth-century political setting and should not be collapsed into a
rule about sharing private customer data.
\href{https://users.manchester.edu/facstaff/ssnaragon/online/texts/318/kant\%2C\%20enlightenment.pdf}{Kant,
\emph{What Is Enlightenment?}}

\Needspace{8\baselineskip}

\hypertarget{M09-self-interest-and-the-burden-of-justification}{%
\section{Self-Interest and the Burden of
Justification}\label{M09-self-interest-and-the-burden-of-justification}}

\noindent \hypertarget{M09-scrutiny-before-reliance}{%
\subsection{Scrutiny Before
Reliance}\label{M09-scrutiny-before-reliance}}

\noindent \noindent A company can gain from a practice while customers bear costs
that the company does not fully internalize. The company may also
control the interface, the outcome definitions, and the evidence
presented to outsiders. These features justify demanding substantiation
before relying on a claim of customer benefit. They do not establish
that every company has committed misconduct. A conflict of interest
identifies a reason to examine a claim and the process that produced it;
a finding of wrongdoing requires evidence about the particular practice.

Hume supplies a strong argument for institutional safeguards. In
\emph{Of the Independency of Parliament}, he treats the assumption of
private-interest pursuit as a principle for designing checks on power.
He explicitly distinguishes that political maxim from a factual claim
that every person is dishonest. The institutional question is whether
the arrangement remains acceptable when participants pursue their
separate interests. Applied to marketing, this asks whether customers
retain meaningful protection when managers respond strongly to sales
targets. The application extends Hume's argument about political
institutions to commercial governance.
\href{https://davidhume.org/texts/emp/ip}{Hume, paragraphs 1--3}

Adam Smith makes the concern more directly commercial. At the end of
Book I, chapter XI of \emph{The Wealth of Nations}, he urges
particularly careful scrutiny of proposals for commercial regulation
made by merchants and manufacturers because their interests can diverge
from the public's. This argument supports examining the incentives and
evidence behind an interested party's proposal. Its original scope is
proposals for laws and regulation, rather than an automatic verdict
against every transaction or product.
\href{https://www.gutenberg.org/files/3300/3300-h/3300-h.htm}{Smith,
Book I, chapter XI, concluding discussion}

The resulting principle for a marketing decision is demanding: when a
firm controls a practice and materially benefits from it, the firm
should substantiate claims of customer benefit with evidence that can be
examined independently. Customers should not have to infer safety or
benefit from the firm's incentives, and an absence of complaints need
not satisfy the burden when complaining is costly or discouraged. This
is an ethical and institutional recommendation, rather than a statement
of a universal legal burden of proof.

\hypertarget{M09-what-would-discharge-the-burden}{%
\subsection{What Would Discharge the
Burden?}\label{M09-what-would-discharge-the-burden}}

\noindent \noindent The burden must be connected to a claim that could be
supported or refuted. For a cancellation process, the firm could supply
evidence on completion, understanding, subsequent charges, and outcomes
for customers who sought to leave. For a personalized recommendation, it
could show how the objective corresponds to customers' stated purposes
and how a person can revise it. For a claim about environmental
performance, it could provide the relevant scope, denominator, and
supporting records. Independent examination need not involve disclosing
identifiable customer records; controlled access and reproducible
aggregates can address different parts of the evidentiary problem.

\begin{figure}[htbp]
\centering
\begin{tikzpicture}
\node[box,text width=3cm] (conflict) at (0,0) {Material conflict\\Firm gains and controls information};
\node[orangebox,text width=3.5cm] (evidence) at (5.3,0) {Independent substantiation\\Customer outcomes and credible comparisons};
\node[box,text width=3cm] (decision) at (10.6,0) {Reasoned decision\\Accept, restrict, redesign, or reject};
\draw[edge] (conflict) -- (evidence);
\draw[accent] (evidence) -- (decision);
\draw[edge,dashed] (decision.south) -- (10.6,-2.5) -| (evidence.south);
\node[align=center,font=\footnotesize] at (8,-2.9) {New evidence can revise the decision};
\node[align=center,text width=12cm] at (5.3,-4.1) {The conflict supplies a reason for scrutiny. The verdict depends on the practice and evidence.};
\end{tikzpicture}

\caption{A conflict of interest changes the scrutiny of evidence; a verdict follows the evaluation of the practice.}\label{fig:M09_burden_justification}
\end{figure}

The sequence separates the reason for scrutiny from the conclusion
reached after scrutiny. A material conflict raises questions about
measurement, selective disclosure, and omitted outcomes. Evidence then
supports acceptance, a narrower claim, redesign, or rejection. A process
that always reaches rejection regardless of the evidence has ceased to
be a critical evaluation.

The standard can also be asymmetric. When the firm receives the gains,
controls the information, and imposes a substantial or
difficult-to-reverse risk on customers, it is reasonable to require
stronger evidence before deployment than before continuing a small,
reversible pilot. That recommendation depends on a stated judgment about
the consequences of acting and waiting. It does not turn uncertainty
into proof of harm. The analyst should identify the information needed,
the feasible alternative, and the condition under which the restriction
can be relaxed.

Consider a retention experiment that reports only revenue. Withhold
acceptance of its customer-benefit claim until the relevant customer
outcomes are supplied. If independent evidence then shows that an easily
refused offer improves informed choice and produces no material
deterioration in the specified outcomes, the original objection must be
revised within that scope. If the additional revenue depends on
preventing an expressed refusal, the evidence supports rejecting that
feature. The recommendation changes because the evidence distinguishes
the mechanism.

\Needspace{8\baselineskip}

\hypertarget{M09-confucian-perspectives-on-marketing-conduct}{%
\section{Confucian Perspectives on Marketing
Conduct}\label{M09-confucian-perspectives-on-marketing-conduct}}

\noindent \hypertarget{M09-righteousness-and-profit}{%
\subsection{Righteousness and
Profit}\label{M09-righteousness-and-profit}}

\noindent \noindent Confucian texts ask what kind of person and relationship a
practice develops, alongside what it produces. In the \emph{Analects}
4.16, the contrast between righteousness and gain concerns the grounds
on which a person judges an action. In 4.5, wealth and status remain
desirable, but their pursuit is constrained by the proper way of
obtaining them. These passages support evaluating the means and
responsibilities of profit seeking. They do not require treating every
profitable exchange as a moral failure.
\href{https://www.gutenberg.org/cache/epub/3330/pg3330-images.html}{Confucius,
\emph{Analects}, Legge translation}

The opening of the \emph{Mencius}, 1A.1, challenges a ruler who frames
the discussion around advantage to his state. Mencius warns that
competition for advantage at every level can damage the relations on
which political order depends. Benevolence and righteousness provide an
alternative basis for judgment. The commercial application is that a
firm cannot settle its responsibilities simply by showing that a
practice raises its own return. The interests of customers, employees,
suppliers, and affected communities require consideration. This extends
a political argument to a contemporary organization.
\href{https://scholarworks.iu.edu/iuswrrest/api/core/bitstreams/265d73a0-6bfa-45df-92ff-4e7d3f8be4b1/content}{Mencius,
1A.1, Eno translation}

For a recommendation system, the question is whether the service helps
the customer pursue a worthwhile purpose while sustaining a trustworthy
relationship. A firm could profit by offering relevant information,
reliable support, and a product the customer continues to value. It
could also increase immediate revenue by exploiting confusion. The fact
that both strategies can be profitable makes the criterion of judgment
consequential.

\hypertarget{M09-trustworthiness-and-the-responsibilities-of-a-role}{%
\subsection{Trustworthiness and the Responsibilities of a
Role}\label{M09-trustworthiness-and-the-responsibilities-of-a-role}}

\noindent \noindent The \emph{Analects} places weight on trustworthiness,
including in its discussion of government in 12.7. Its account of the
rectification of names in 13.3 concerns the relation between language,
social roles, and proper conduct. A role is associated with standards of
performance. Calling someone a trustworthy adviser does not supply the
conduct that warrants that description. These ideas can inform a
contemporary examination of how marketing labels create expectations.
\href{https://scholarworks.iu.edu/bitstreams/d1f62d31-9082-41d0-999f-4c3efa0c8857/download}{Confucius,
\emph{Analects}, 12.7 and 13.3; Eno translation}

A platform that calls a paid placement a recommendation should examine
what customers understand that description to mean. If users infer that
the ranking serves their goals while the rule principally favors an
undisclosed payment, the claimed role and actual conduct diverge.
Likewise, a customer-care team whose only rewarded outcome is preventing
cancellation may face incentives inconsistent with the service its name
promises. Clear labeling is one response; aligning responsibilities,
rewards, and actual conduct addresses the underlying relation.

Confucian role ethics does not simply duplicate Kant's account of
autonomy. It emphasizes cultivated dispositions, relationships, and
responsibilities within social roles. Historical Confucian roles also
contain hierarchies that require critical examination when applied to
contemporary customers and organizations. A modern application should
state why the proposed responsibility is justified and how affected
people can challenge its interpretation. Calling a manager benevolent
cannot make unilateral control over a customer's goals self-justifying.

\hypertarget{M09-remonstrance-reciprocity-and-learning}{%
\subsection{Remonstrance, Reciprocity, and
Learning}\label{M09-remonstrance-reciprocity-and-learning}}

\noindent \noindent The \emph{Analects} 13.23 distinguishes harmony from mere
agreement, and 14.22 connects service to a ruler with truthful
challenge. Applied to a marketing team, these ideas support a
responsibility to raise a material objection even when the proposal is
commercially attractive. Harmony need not require an analyst to suppress
evidence that the team's preferred metric conceals customer harm. The
organizational question is whether disagreement can be voiced and
considered without being treated as disloyalty.

Reciprocity in \emph{Analects} 15.24 encourages examination of how one
treats others. In marketing, it can motivate asking whether an offer,
disclosure, or cancellation process is one the decision maker could
reasonably accept from the customer's position. Different customers have
different resources and purposes, so the manager's own tolerance for
inconvenience is insufficient. Perspective taking should lead to inquiry
about affected people, including those who face language, accessibility,
or financial constraints.

The conjunction of learning and reflection in \emph{Analects} 2.15 also
fits AI-assisted analysis. Accumulating generated explanations without
examining them leaves the learner dependent on assertions. Reflection
without engaging evidence can become speculation. A useful learning
exercise requires the student to compare a generated argument with the
source, inspect its premises, and revise a conclusion when the evidence
warrants it. These are contemporary applications of the passages, rather
than claims that Confucius anticipated automated marketing systems.

\hypertarget{M09-mencius-and-xunzi-on-cultivating-conduct}{%
\subsection{Mencius and Xunzi on Cultivating
Conduct}\label{M09-mencius-and-xunzi-on-cultivating-conduct}}

\noindent \noindent Confucianism contains substantial internal disagreement.
Mencius's account of moral beginnings, including compassion in 2A.6,
emphasizes capacities that can be developed. Xunzi's chapter on human
nature emphasizes the disorder that can follow unchecked desires and the
importance of learning, ritual, and deliberate effort in producing good
conduct. The differing conceptions of human nature should not be reduced
to a single slogan about whether all people are good or bad.
\href{https://scholarworks.iu.edu/iuswrrest/api/core/bitstreams/265d73a0-6bfa-45df-92ff-4e7d3f8be4b1/content}{Mencius,
2A.6};
\href{https://afe.easia.columbia.edu/main_pop/ps/ps_china-xunzi-human-nature.htm}{Xunzi,
chapter 23, Columbia selection};
\href{https://plato.stanford.edu/entries/xunzi/}{Goldin, \emph{Xunzi}}

Xunzi offers an especially useful response to reliance on goodwill
alone. Ethical conduct requires sustained formation and practices that
direct behavior. In a firm, a contemporary application would examine
what employees repeatedly learn to notice, what the reward system
encourages, and how the organization handles a truthful report of an
unwanted outcome. Ritual propriety, often rendered \emph{li}, concerns
patterned social conduct and cultivation; it should not be equated
mechanically with a compliance checklist. The word for profit or
benefit, also commonly romanized \emph{li}, is a different Chinese word.

For the subscription case, a Confucian evaluation can examine whether
the customer-care role is performed faithfully, whether a refusal is
treated with respect, and whether the company sustains a relationship
worthy of trust. A Kantian evaluation emphasizes the principle of action
and treatment of the customer as an end. Hume's institutional argument
asks how the process behaves under strong private incentives. The
perspectives can support the same redesign for different reasons.
Keeping those reasons explicit permits disagreement and avoids
presenting distinct philosophical traditions as interchangeable.

\Needspace{8\baselineskip}

\hypertarget{M09-decisions-that-can-withstand-criticism}{%
\section{Decisions That Can Withstand
Criticism}\label{M09-decisions-that-can-withstand-criticism}}

\noindent \noindent A decision memorandum should begin with the action and the
reason for it. For the subscription case, a supportable recommendation
could retain an optional offer while making cancellation immediately
effective when requested. Its justification would distinguish the
measured commercial effect, the customer's stated objective, the
relevant ethical principle, and the applicable legal requirements. If
the firm's gain depends on preventing an understood refusal, increased
revenue cannot by itself answer the ethical objection.

Consider three further decisions. A retailer wants to display a
personalized discount based on past purchases. Assess whether the offer
changes the feasible options, conceals material terms, or interferes
with refusal. A product team wants to replace a nutrition objective with
a click objective because clicks are easier to optimize. Identify the
customer purpose and the missing connection between that purpose and the
metric. An advertiser wants to describe one improvement as evidence of a
company-wide transition. Specify the population, denominator, and time
horizon needed to support the broader claim.

For each decision, formulate the strongest defensible argument in favor
of the practice, the most consequential objection, and the premise that
separates them. State what evidence would alter the recommendation. A
critical conclusion can endorse a practice, narrow its claim, improve
its design, or reject it. Its strength comes from the reasons and
evidence that survive examination.

\Needspace{8\baselineskip}

\hypertarget{M09-reading-path}{%
\section{Reading Path}\label{M09-reading-path}}

\noindent \noindent Begin with Kant's introduction to the Transcendental Dialectic
and the \emph{Prolegomena}, §§50--54. Read the third antinomy alongside
the distinction between appearances and things in themselves. The
appendix on the regulative use of ideas explains the constructive role
of reason. Grier's
\href{https://plato.stanford.edu/entries/kant-metaphysics/}{account of
Kant's critique of metaphysics} provides an interpretive guide. Read the
\emph{Groundwork}, second section, separately for the universal-law and
humanity formulations; these supply the ethical analysis rather than a
conclusion derived from the Dialectic alone.

The \href{../01_Resources/15_Marketing_Dialectics/README.md}{categorized
research collection} links the primary texts, scholarly interpretations,
marketing research, and industry cases. It identifies full texts,
excerpts, versions, and remaining access gaps. The two ASA decisions,
the advertising-measurement study, and the subscription-interface
research supply contrasting empirical and institutional settings for
assessing the scope of a marketing argument.

For scrutiny of interested claims, read Hume's opening paragraphs on
institutional checks and Smith's conclusion to Book I, chapter XI. For
Confucian perspectives, compare \emph{Analects} 2.15, 4.5, 4.16, 12.7,
13.3, 13.23, 14.22, and 15.24 with \emph{Mencius} 1A.1 and 2A.6 and
Xunzi's discussion of human nature. The resource guide distinguishes the
original settings, translation choices, and proposed marketing
applications.

\end{document}
